\section{The Intellectual Soul}

\label{sec:IntellectualSoul}

The Western Tradition teaches the four Kingdoms and that Man participates in all of them through body and soul. In common with other traditions, there are different layers to the soul, as shown in diagram \ref{tab:LayersSoul} adapted from \textbf{St Thomas Aquinas}.

\begin{table}[h]\centering
\begin{tabular}{rccc}\toprule\small
  &
\textbf{Kingdom} &
\textbf{Hermetic} &
\textbf{Source}\\\midrule
Body &
Mineral &
Body &
Natural\\
Nutritive Soul &
Vegetable &
Etheric Body &
Natural\\
Sensitive Soul &
Animal &
Astral Body &
Natural\\
Intellectual Soul &
Man &
Causal Body &
Supernatural\\\bottomrule
\end{tabular}
\caption{Layers of the Soul}
\label{tab:LayersSoul}
\end{table}
In man, these form a unity, unlike some gnostic and theosophical systems which try to separate them. As Aquinas puts it, ``the sensitive soul, the intellectual, and the nutritive soul are in a man, and are numerically one and the same soul." But the body, also, is a part of man. Unlike dualist systems which view the soul and body as separate substances, or monist systems such as materialism or spiritualism which deny the existence of the soul or body, respectively, this system is nondual. That is, the soul as spirit is the form of the body as matter.

This is the traditional teaching, which regards the material world as the reflection of the spiritual. Theologians have not really drawn out fully the consequences of this teaching. In particular, it means that all aspects of the sensitive soul, nutritive soul, and body take on the form of the intellectual soul. This has to include our racial and physical characteristics as well as our mental characteristics.

Aquinas makes this clear in his doctrine of subordination, which is accepted by both Evola and Guenon. He writes:

\begin{quotex}
the sensitive is subordinate to the intellective and the nutritive to the sensitive, as potency is subordinate to act, since in the order of generations the intellective comes after the sensitive and the sensitive after the nutritive. 

\end{quotex}
Yet things are not so simple. Aquinas continues:

\begin{quotex}
man's intellectual soul contains virtually whatever belongs to the sensitive soul of brute animals, and to the nutritive soul of plants. 

\end{quotex}
So although the lower parts of the soul and the body are passive in relation to the intellectual soul, this is only virtual, not actual. That is why we see man as fragmented. Instead of being guided by his intelligence, his acts are motivated by physical desires, conflicted emotions, and personal whims. Man comes to see this as his ``normal" and ``natural" condition. If he hears otherwise, he will dismiss it as being ``life denying", prudish, authoritarian, uptight, and so on. Such a person, similarly to Nietzsche, will regard intelligence as merely instrumental; that is, it will help him devise the means to satisfy his desires. However, he does not use his intelligence to determine the ends themselves.

Even when a man is drawn to such a teaching by virtue of his spiritual condition or destiny, in today's world there is nowhere he can turn. The philosophy books explain all this in the abstract, while what he really needs is a direct intuitive understanding or gnosis. In the Hermetic schools, he would learn to observe the various movements of his soul just as the physicist observes the movements of physical bodies. In this way he learns to distinguish between the processes of the various parts of the soul. He can catalog them and, through various exercises, develop the skills to actualize the proper relationships among those parts.

As he learns to see things in the clear light of the intellect, he will develop a sense of detachment and will no longer be driven by the irrational desires of the body nor be perturbed by emotional swings. This is not to say that he will not experience higher emotions such as equanimity and peacefulness, just that they are not the motivating factor. Only when a man acts from the intellectual soul is he free. Then he will possess the virtues: \textbf{temperance}, \textbf{courage}, \textbf{prudence}, and \textbf{justice}.



\flrightit{Posted on 2011-11-14 by Cologero }

\begin{center}* * *\end{center}

\begin{footnotesize}\begin{sffamily}



\texttt{Gabe Ruth on 2011-11-16 at 10:06 said: }

I've been lurking for a few weeks, trying to figure out what's going on here. Some things have troubled me, but this take on Nietzsche convinces me that you are very right about many things, and quiets my fears about your purpose. 

I think the apostasy of a man like Nietzsche is one of the strongest pieces of evidence in the indictment of the exoteric church. Though I think the zeitgeist bears most of the responsibility, surely a more vigorous, confident church would have been able to reach a man with such great incite into the human condition. When I read accounts of his decline, I was reminded of Weston's ravings to Ransom in Perelandra, when he regains control temporarily.

Given your view of the modern church, I wonder what you think of the Theology of the Body. While I doubt you think very well of JP2, those ideas seem to be a fleshing out of some of what you get into here, and not without value.


\hfill

\texttt{charlotte on 2011-11-16 at 16:33 said: }

JPII actually restored the Latin Mass in some places, but was unable to get the method of priestly ordination changed. Recently I've been thinking a lot about the assassination of JPI, the prophecies of Fatima and so on. Benedict is making me wonder, as well, with all this stuff about a global banking system, in which the Vatican would no doubt play a significant role…All in all I'm very inclined to go seek out one of those traditional churches. there are a few places in England that still offer a traditional high mass and I feel a great need to go find one as soon as possible.


\hfill

\texttt{Cologero on 2011-11-17 at 12:20 said: }

I am not in a position to comment on the Theology of the Body. Karol Wojty?a was part of a movement that combined Thomism with Phenomenology. In principle, that sounds fruitful since it seems to make the metaphysical experiential. I don't know where that led to or what part it plays. But it is the necessary approach.

This is not something we can do on a blog, just to make people aware of that option. For example, as for the Traditional teaching on the various souls … is that something theoretical or experiential? Was Aquinas just weaving words or was he describing an inner experience? Disputing and debating positions goes on forever with no result.

So, as far as a theology of the body … we know the proper relationship is spirit –{\textgreater} soul –{\textgreater} body. However, we find that the body does not simply respond to the will. Why is that? How can that change? And so on. So the approach is not a mental exercise, but rather a spiritual exercise.


\hfill

\texttt{Graham on 2011-11-18 at 10:37 said: }

``In particular, it means that all aspects of the sensitive soul, nutritive soul, and body take on the form of the intellectual soul. This has to include our racial and physical characteristics as well as our mental characteristics."

The Greeks held that an attractive body is a sign of interior worth; this is particularly revealed by the word `kallis' which could mean both physical beauty and goodness. One finds, curiously, that Plato believed looks should be considered (to a small degree) in the selection of the golden-souled guardians in his perfect republic.

It is an instinctively attractive proposition – and self-examination as well as scientific studies bear out that we think like that to this day – but we also know from experience that it often doesn't hold true. How does Gornahoor's `theology of the body' navigate this issue? Does the good-looking person have a greater potential for goodness?


\hfill

\texttt{Gabe Ruth on 2011-11-18 at 11:55 said: }

From Cologero's comment: ``the body does not respond to the will". This is referring to the body's urges, but I think physical appearance would be even less connected to the will of the spirit.

The spirit takes precedence, as the immortal component. The fact of aging also indicates that physical appearance is not an exact indication of spiritual condition. 

On the other hand, who can say how completely the modern conception of who is good-looking has been distorted?

Regarding the Theology of the Body, I am far from an authority and have never even read the primary text. But making the metaphysical experiential is very much the goal. It is an exoteric teaching, but a metaphysically ambitious one that contemplates what human nature as male and female says about the nature of God. It also seeks to ground Catholic sexual ethics in metaphysics as well.

Charlotte, what was he trying to change about the method of priestly ordination?

I just read an older post on the Aztecs, and it brought the movie ``The Mission" to mind. What are your thoughts on the Jesuit Reducciones?


\hfill

\texttt{Graham on 2011-11-18 at 13:15 said: }

Gabe,

If the doctrine of correspondences is to have any validity, the physical body must be a reflection of the spirit. The post spells it out clearly, saying that the soul as spirit is the form of the body as man, and that this applies to physical characteristics as well as mental. So the problem remains.

In this view, growth and aging would have to do with the `phasing' of the intellectual soul into and out of the material world. The prime of life would be the fullest physical and mental expression of the spiritual form, and it's for this reason that Aquinas asserts that the resurrected body will appear as that of a 28 year old. So aging does not alter the problem.

As for your final question, I agree that the modern American taste is different from that of the ancient Greek, and I go further in asserting that the taste of the dispassionate Plato differed again from that of his average carnal contemporary. His sense of human beauty would have been subtle and spiritualized. However, since number is eternal, and beauty has its basis in proportion (or vice versa), there are also unchanging givens at play. A proportionate face is a proportionate face: there would be significant overlap between all these differing tastes. More concretely, we know from Greek sculpture that their idea of human beauty was much like ours, if more sublime.

It is odd though that Plato would put ideas like this in the mouth of the notoriously ugly Socrates. It seems like one of those contradictions that Strauss (speaking of him) would want to explore further.


\hfill

\texttt{charlotte on 2011-11-19 at 05:20 said: }

Jesus Christ was/is extremely beautiful….

Grail Queens tend to be also, but these are ultimate examples.

More generally, if somebody is physically very attractive (as opposed to just very pretty) it indicates a certain charismatic power of radiation, and love energy radiates outwards. In contrast with someone who is simply good looking but `fake' or very vain, for instance, which is a contracting force. Those people are often appealing at first sight but soon lose magic and mystery unless the eye of the beholder is similarly vain. You often find that enormous sexual power equates to spiritual power in some degree, because it's a measure of that person's overall energy (as with the `man of action) – we have 7 centres and the same force flows through them all. However, if it `gets stuck' in the sex centre then that's not `good', although the potential is there. Throughout one's life I would imagine that part of the developmental task is to keep the energy flowing smoothly and constantly. 

Understanding how to keep this energy at peak condition, to master it rather than allowing it to master us, is very important and involves a lot of polarity working (which is why the sexual and soul-mate relationships are so important during our lives). Highly spiritualised people will have their energy primarily emanating from the heart centre, while the crown will be open,shining and receiving force from the vertical hierarchy. Other people will be naturally inclined to love them – to love the angelic, Christed spirit flowing through the individual. Buddhist training helps us to master the energy flow and strengthen the centres.

Regarding the priestly ordination, I understand that during Vatican II one of the changes made to the church was with respect to this. JPII was unable to revert back to the original form but did win some victories with respect to the traditional mass. perhaps this is part of the mysterious spiritualisation process of the church, because it is still possible to be part of a catholic community, even if it is not necessarily visible…


\hfill

\texttt{Cologero on 2011-11-19 at 16:35 said: }

Regarding physical attractiveness and related issues, it would be better to avoid speculation in favour of what we can directly know.

First of all, the body is a means of self-expression. We smile, frown, show anger, excitement, and so on. The body then reveals our inner states. Then, it is a means of our self-determination. We have projects, aims, goals and we execute them through our bodies. Thus, it is a means of our creative self-expression.

We then note our privations. We may spontaneously express anger, seemingly against our will. We may stammer in am embarrassing situation. We want to exercise but our body resists. These are all things to observe in developing the will.

Keep in mind that we are born into a particular stream of soul-life and material conditions, so it can't be said the soul and body are a perfect expression of the spirit ab initio. This is also a task, not just a given.


\hfill

\texttt{charlotte on 2011-11-20 at 03:33 said: }

well I think whichever way one looks at it beauty is a form of power, and like any form of power can be used for good or ill. Of course beauty of spirit radiates outwards and makes a person attractive in its own right, certainly lovable, regardless of whether they have perfect bodies and flawless skins.


\end{sffamily}\end{footnotesize}
